\chapter{What the bowl remembers}\label{connection:ferments}
\regionline{Doughs, batters and a portion kept back}
\atlasfiguretitle{What the bowl remembers: paired comparison}
\comparisonplate{Bug}{Injera}{A culture kept in a jar and a fermented bread made from batter.}

A starter makes the next loaf begin before the flour is weighed. A little mixture has been kept, fed or revived, and the baker expects it to do useful work again. That expectation links breads which may look almost nothing alike. A Finnish rye round, an Ethiopian or Eritrean injera and Māori rēwena can all begin with a culture entrusted with the next batch.

The Finnish word juuri, root, gives the kept leaven a wonderfully material name. In the ruisleipä story, Eric Bergman recovers his grandmother's frozen dough and learns to bake with it. Later, he dries dough for his move to Zagreb and revives it there. The thing carried is small, but it makes a familiar bread possible in a new kitchen.\citep{finroot,finjourney}

\section{A living mixture has a maker}
Whiti Hereaka's rēwena demonstration in Te Ara shows a potato-based bug mixed with flour and water, left to become active and then worked into bread. The sequence is photographed through a named maker's hands. It gives a reader a way to see the culture's activity before encountering a loaf.\citep{rewena}

The practical rēwena recipe in this book follows a separately credited formula and includes a starter build. Its jar deserves attention because the bake depends on reliable activity. Look at the rise after a feed, the bubbles through the mixture and what happens when that portion enters the final dough. The word sourdough opens a useful comparison, but it need not flatten a potato culture and a rye root into the same household practice.

Injera takes the comparison from dough to batter. The retained ferment, ersho, begins another batch, while a cooked portion of batter, absit, becomes part of the preparation before cooking. The baker watches fermentation and flowing consistency rather than looking for a kneaded ball. An active culture matters, yet the mixture must also spread across its cooking surface and set into a bread that can be lifted.\citep{maskalstarter,maskalinjera}

\section{The mixture changes the questions}
Kisra's sorghum batter asks for a spreading action. Appam's rice-based batter asks for a pan that shapes a thin edge and deeper centre. With each, the cook can observe how bubbles, flow and heat meet in a small test bread. The useful question is often whether the mixture will make the intended surface, rather than whether it behaves like wheat dough.\citep{moggakisra,familiarappam06}

This is why a comparison across grains should preserve the recipe's vocabulary. Stretching a wheat dough, pouring an injera batter and feeding a rēwena bug are different actions. Their relationship lies in how cooks learn to read a mixture that changes with time. A clock helps arrange the day, but the bowl tells the baker whether that day is proceeding as expected.

The home recipes include explicit starter proportions because a reader needs an amount that can be weighed. They also include readiness cues because a living culture and a warm room do not always agree with a printed schedule. Record both. A note that says how the mixture looked is more useful next time than a duration copied from a book.

\section{Time can be arranged in other ways}
A poolish lets the baguette baker ferment part of the flour with commercial yeast before mixing the final dough. Campagne uses a sourdough culture. Soda bread and the modern damper printed here give the cook a different timetable, with chemical raising agents and a prompt move to the oven. These are different ways to make bread fit the day, rather than positions on a ladder of seriousness.\citep{baguetterecipe,campagnerecipe,samdamper09}

The bread that can be made now has its own value. A pan of pol roti, the warm stack of Fijian roti or a quick soda loaf can bring fresh bread to a meal without maintaining a culture. The cook's available time is part of the recipe's setting. An atlas becomes more welcoming when it shows that pleasure is possible through several rhythms.

Shokupan introduces another kind of advance work. Its yudane scald is prepared ahead, but scalding flour is not fermentation. The cooked paste used for the Agege-style recipe likewise needs to be distinguished from a sourdough starter. A bowl set aside overnight can contain several kinds of preparation, and the recipe tells us what that waiting is for.\citep{urashokupan07,woodardagege}

\section{Keeping more than an ingredient}
A culture can also carry a method: when it is fed, how much is retained and which signs a particular baker trusts. That knowledge is easier to lose than a jar is to replace. When someone gives you a starter, ask how they use it and write down the answer beside your first bake. The notes can hold the connection between a person, a portion of mixture and a bread you hope to make again.

\section{A rise beneath a lid}
Mantou offers another useful turn in the comparison. Its printed recipe uses commercial yeast, yet the rising dough is prepared for a steamer rather than a browned oven crust. The smooth surface, space between portions and steady cooking steam matter to the finished bun. Fermentation is one part of its preparation, while the cooking environment gives that risen dough another form.\citep{elainemantou07}

Wei Guo's memory of her parents preparing dozens of mantou before Chinese New Year follows the dough into family work and festive shapes. Plain portions, flowers and animals can begin within the same breadmaking practice. A living mixture becomes something to recognise as well as something to eat.\citep{guomantou07}

\begin{breadnote}{Follow the portion kept back}
Read ruisleipä on page~43, injera on page~83 and rēwena on page~167. Choose one culture to begin with. Keep its feeding method distinct from the others, and record activity and temperature before making the main recipe.
\end{breadnote}
